\begin{pblock}{From topography to habitat}
The DEM is an intermediate product. Three layers follow from it, each kept
honest about what it does not know:

\begin{itemize}
  \item \term{Hydroperiod} — the fraction of \emph{time} a pixel is
        submerged, evaluated over a continuous modelled tide series and
        deliberately not over the satellite dates, which inherit the orbit's
        sampling bias. This, not water frequency, is what ecology needs.
  \item \term{Hypsometry} — the area--elevation curve. Convex means accreting,
        concave sediment-starved, and it gives the flooded-area response to
        sea-level rise directly.
  \item \term{Change between epochs} — never a bare difference of two DEMs.
        Each carries a significance mask from the propagated uncertainty
        $\sigma_\Delta=\sqrt{\sigma_1^2+\sigma_2^2}$, because a datum shift
        between epochs otherwise masquerades perfectly as erosion.
\end{itemize}
\end{pblock}
